\chapter{Metodología}
\label{cap:metodologia}

Este capítulo describe el enfoque metodológico seguido para alcanzar los
objetivos planteados, las herramientas y entornos empleados, y los criterios de
evaluación previstos para valorar el sistema. La descripción detallada del
diseño y la implementación se reserva para el Capítulo~\ref{cap:desarrollo}.

%% =============================================================================
\section{Metodología de desarrollo}
\label{sec:metodologia-desarrollo}

El trabajo se ha abordado mediante un enfoque de desarrollo \textbf{iterativo e
incremental}, organizado en sprints sucesivos cuya gestión se ha apoyado en un
tablero \textit{Kanban}. En este tablero, las tareas avanzan por columnas según
su estado, lo que ha permitido mantener una visión clara del progreso y priorizar
el trabajo pendiente en cada iteración. En lugar de construir la solución completa
de una sola vez, se ha avanzado por incrementos funcionales, validando cada parte
antes de integrarla con el resto.

El orden de construcción ha seguido la cadena de dependencias del sistema. En
primer lugar se levantó la infraestructura de datos (las bases de datos, el bus de
mensajería y los servicios de almacenamiento y visualización), por ser el soporte
sobre el que se apoya todo lo demás. A continuación se desarrolló el backend, que
centraliza la lógica de negocio y la \gls{api}. Después se abordó la capa robótica
(\texttt{robot\_ws}), que produce los datos y consume los servicios del backend. Y,
por último, se implementó la aplicación de usuario, que también consume la
\gls{api} del backend.

La integración entre el robot y el backend se realizó funcionalidad a
funcionalidad, siguiendo siempre el mismo patrón: primero se implementaba en el
backend el \textit{endpoint} o la lógica necesaria para una capacidad concreta, y
a continuación se desarrollaba la parte del robot que se conecta a ella. Este
encadenamiento, repetido para cada funcionalidad (telemetría, incidencias,
comandos de control y autenticación), garantizaba que cada extremo dispusiera ya
de su contraparte y que cada incremento fuera verificable de principio a fin.

Finalmente, todo el código se ha organizado como un \textbf{monorepo}, es decir,
un único repositorio que agrupa todos los componentes del sistema, lo que facilita
la coherencia entre las distintas partes y simplifica el despliegue conjunto. El
repositorio se estructura en cuatro bloques principales: la capa robótica
(\texttt{robot\_ws}), el backend (\texttt{backend}), la aplicación de usuario
(\texttt{flutter\_frontend}) y un conjunto de \textit{scripts} de automatización
del despliegue. La arquitectura concreta de cada bloque se detalla en el
Capítulo~\ref{cap:desarrollo}.

%% =============================================================================
\section{Herramientas, frameworks y entornos}
\label{sec:herramientas}

El sistema se apoya en un conjunto de tecnologías de código abierto, seleccionadas
por su madurez y por su integración con el ecosistema de \gls{ros}. La
Tabla~\ref{tab:tecnologias} resume las principales herramientas empleadas en cada
capa del sistema.

\begin{table}[htbp]
  \centering
  \caption{Principales tecnologías empleadas por capa del sistema. Fuente:
  elaboración propia.}
  \label{tab:tecnologias}
  \begin{tabular}{@{}lp{8.5cm}@{}}
    \toprule
    \textbf{Capa} & \textbf{Tecnologías} \\
    \midrule
    Plataforma robótica & TurtleBot~4, \gls{ros} Humble \\
    Navegación y mapeado & \gls{nav2}, slam\_toolbox, explore\_lite \\
    Percepción & \gls{yolo}v8n (Ultralytics), OpenCV \\
    Mensajería & Apache Kafka (cliente confluent-kafka) \\
    Backend & FastAPI, SQLAlchemy, Uvicorn \\
    Persistencia & PostgreSQL, InfluxDB, MinIO \\
    Visualización & Grafana \\
    Aplicación de usuario & Flutter \\
    Infraestructura & Docker, Docker Compose \\
    \bottomrule
  \end{tabular}
\end{table}

A continuación se detalla el papel de cada tecnología y la justificación de su
elección, agrupadas por capa.

\subsection{Plataforma robótica y navegación}

El robot empleado es el TurtleBot~4, una plataforma de robótica móvil de
referencia compuesta por una base \textit{Create~3}, un computador de a bordo
(Raspberry~Pi~4), un sensor \gls{lidar} y una cámara OAK-D. Se ha elegido por ser
una plataforma estándar, asequible y con soporte oficial para \gls{ros}, lo que
evita desviar el esfuerzo del trabajo hacia la integración de hardware a bajo
nivel. Sobre ella se ejecuta \gls{ros} en su distribución \textit{Humble}, una
versión de soporte prolongado que es la recomendada para el TurtleBot~4
\parencite{macenski2022ros2}.

La navegación se construye sobre tres paquetes consolidados del ecosistema. La
pila \gls{nav2} se ocupa de la planificación y el control del movimiento
\parencite{macenski2020nav2}; \texttt{slam\_toolbox} construye el mapa del entorno
mediante \gls{slam} láser \parencite{macenski2021slamtoolbox}; y
\texttt{explore\_lite} dirige la exploración autónoma inicial que recorre el
espacio desconocido hasta completar el mapa. Reutilizar estas soluciones, en lugar
de reimplementar la navegación, permite concentrar el trabajo en la orquestación y
la integración del sistema.

\subsection{Percepción visual}

La detección de personas se realiza con \gls{yolo} en su versión \textit{nano}
(YOLOv8n), a través de la biblioteca de Ultralytics
\parencite{jocher2023ultralytics}. Se ha optado por la variante \textit{nano} por
ser la más ligera de la familia, lo que permite ejecutar la inferencia en la
\gls{cpu} del propio robot a una velocidad suficiente para la patrulla, sin
necesidad de una \gls{gpu} dedicada. La biblioteca aporta, además, modelos
preentrenados y una interfaz sencilla que reduce el coste de integración. El
tratamiento de las imágenes, esto es, la captura de los fotogramas de la cámara y
la codificación del clip de vídeo de cada incidencia, se apoya en OpenCV,
biblioteca estándar de facto en visión por computador.

\subsection{Mensajería orientada a eventos}

El bus de eventos que comunica el robot con el servidor es Apache Kafka
\parencite{kreps2011kafka}. Se ha preferido a una solución de mensajería más
ligera porque conserva los mensajes en un registro durable, lo que desacopla por
completo al productor del consumidor y permite que el backend procese la
telemetría y las incidencias a su propio ritmo, e incluso reprocese eventos
pasados. La comunicación se organiza en tres \textit{topics}
(\texttt{robot\_telemetry}, \texttt{robot\_incidents} y \texttt{robot\_commands}).
Tanto el robot como el backend acceden a Kafka mediante \texttt{confluent-kafka},
un cliente de alto rendimiento basado en la biblioteca \textit{librdkafka},
escrita en C.

\subsection{Backend y servicios web}

El backend se ha desarrollado con FastAPI, un framework web asíncrono para Python.
Su elección responde a tres motivos: el alto rendimiento que aporta su modelo de
servidor asíncrono (\textit{ASGI}), la validación automática de los datos de
entrada y salida mediante Pydantic, y la generación automática de la documentación
de la \gls{api} según el estándar OpenAPI. La aplicación se sirve con Uvicorn, un
servidor \textit{ASGI}, y la configuración (credenciales y direcciones de los
servicios) se gestiona con \texttt{pydantic-settings} a partir de variables de
entorno, lo que facilita el despliegue en distintos equipos.

El acceso a la base de datos relacional se realiza mediante el \gls{orm}
SQLAlchemy, que abstrae las consultas \gls{sql} y mapea las tablas a clases de
Python, mientras que Alembic versiona y aplica los cambios del esquema. El
almacenamiento de los vídeos se gestiona con \texttt{boto3}, el \gls{sdk} de
\gls{s3}, que permite operar sobre MinIO con la misma interfaz que el servicio de
Amazon. Por último, la seguridad se apoya en \texttt{pyjwt} para los tokens
\gls{jwt} y en \texttt{passlib} con \textit{bcrypt} para el cifrado de las
contraseñas de los usuarios.

\subsection{Persistencia, almacenamiento y visualización}

La distinta naturaleza de los datos del sistema ha llevado a emplear un almacén
específico para cada tipo. Los datos estructurados y con relaciones (usuarios,
robots e incidencias) residen en PostgreSQL, una base de datos relacional madura
que garantiza la integridad mediante transacciones. La telemetría del robot
(batería, posición y estado), de carácter temporal y alta frecuencia, se almacena
en InfluxDB, una base de datos de series temporales (\gls{tsdb}) optimizada para
este tipo de datos y para las consultas por intervalos de tiempo. El vídeo de las
incidencias, al ser un objeto binario de gran tamaño, se guarda en MinIO, un
almacén de objetos compatible con la \gls{api} de \gls{s3}; esta separación evita
sobrecargar la base de datos relacional con ficheros pesados. Finalmente, Grafana
proporciona la visualización de la telemetría mediante paneles que consultan
InfluxDB, con un panel por robot que la aplicación de usuario incrusta
directamente.

\subsection{Aplicación de usuario}

La aplicación con la que interactúa el usuario se ha construido con Flutter, el
framework de Google basado en el lenguaje Dart. Su principal ventaja es permitir
desarrollar, a partir de una única base de código, una aplicación para móvil, web
y escritorio, algo idóneo para un trabajo con tiempo limitado. Sobre Flutter se
apoyan varios paquetes especializados: la gestión del estado se resuelve con
\texttt{flutter\_riverpod}; la navegación entre pantallas, con \texttt{go\_router};
las llamadas a la \gls{api} \gls{rest} del backend, con el cliente HTTP
\texttt{dio}; la persistencia local de la sesión, con \texttt{shared\_preferences};
los componentes de la interfaz, con \texttt{shadcn\_ui}; y la apertura del panel de
Grafana en el navegador, con \texttt{url\_launcher}.

\subsection{Infraestructura y despliegue}

Toda la infraestructura de datos (PostgreSQL, Kafka, InfluxDB, MinIO y Grafana) se
ejecuta en contenedores gestionados con Docker y orquestados mediante Docker
Compose. Esta decisión persigue la reproducibilidad del entorno: un único fichero
declarativo describe todos los servicios, sus versiones y su configuración, de modo
que el sistema completo puede levantarse en cualquier equipo con un solo comando,
sin instalaciones manuales y con un comportamiento idéntico. El aislamiento entre
contenedores evita, además, los conflictos de dependencias entre los distintos
servicios.

\subsection{Entorno de desarrollo y herramientas de apoyo}

El grueso del sistema se ha desarrollado en \textbf{Python}, lenguaje común al
backend y a los nodos del robot, lo que ha permitido compartir conocimientos y
bibliotecas entre ambas partes. Es, además, el lenguaje de referencia para la
programación de nodos en \gls{ros} y para el ecosistema de visión por computador y
aprendizaje automático empleado en la percepción. Para aislar las dependencias de
cada componente se han utilizado entornos virtuales (\textit{virtual environments},
\texttt{venv}), de modo que el backend y la capa robótica mantienen conjuntos de
paquetes independientes y sin interferencias entre versiones; las dependencias de
cada entorno se declaran de forma explícita (por ejemplo, en el fichero
\texttt{requirements.txt} del backend), lo que refuerza la reproducibilidad ya
mencionada.

Durante la elaboración del trabajo se ha empleado, asimismo, el asistente de
inteligencia artificial Claude como herramienta de apoyo en tareas de revisión,
búsqueda de información (\textit{research}) y documentación. Su uso se ha limitado
en todo momento a un papel auxiliar: las decisiones de diseño del
sistema han sido tomadas y validadas por el autor, sin que la herramienta las
sustituya.

\subsection{Bibliotecas y dependencias principales}

A modo de síntesis, la Tabla~\ref{tab:librerias} recoge las bibliotecas más
relevantes empleadas en el desarrollo, indicando su función y el componente del
sistema en el que se utilizan. La mayoría se han presentado ya en las subsecciones
anteriores; la tabla las reúne como referencia.

\begin{table}[htbp]
  \centering
  \small
  \caption{Bibliotecas y dependencias principales, su función y el componente en
  el que se emplean. Fuente: elaboración propia.}
  \label{tab:librerias}
  \begin{tabular}{@{}lp{9cm}@{}}
    \toprule
    \textbf{Biblioteca} & \textbf{Función} \\
    \midrule
    \multicolumn{2}{@{}l}{\textit{Backend (Python)}} \\
    FastAPI & Framework web asíncrono y exposición de la \gls{api} \gls{rest} \\
    Uvicorn & Servidor \textit{ASGI} que ejecuta la aplicación \\
    SQLAlchemy & \gls{orm} de acceso a PostgreSQL \\
    Alembic & Migraciones del esquema de la base de datos \\
    \texttt{influxdb-client} & Escritura de la telemetría en InfluxDB \\
    \texttt{confluent-kafka} & Consumo de los eventos del bus Kafka \\
    \texttt{boto3} & Acceso a MinIO (\gls{s3}) para el vídeo de incidencias \\
    Pydantic & Validación de datos y configuración por entorno \\
    PyJWT & Emisión y verificación de tokens \gls{jwt} \\
    \texttt{passlib} / \textit{bcrypt} & Cifrado de las contraseñas de usuario \\
    \texttt{httpx} & Cliente HTTP asíncrono \\
    \midrule
    \multicolumn{2}{@{}l}{\textit{Capa robótica (\texttt{robot\_ws}, Python sobre \gls{ros})}} \\
    Ultralytics (\gls{yolo}) & Detección de personas durante la patrulla \\
    OpenCV & Captura de fotogramas y codificación del clip de vídeo \\
    NumPy & Cálculo numérico de soporte a la percepción \\
    \texttt{confluent-kafka} & Publicación de telemetría e incidencias en Kafka \\
    \texttt{boto3} & Subida de los clips de vídeo a MinIO \\
    \texttt{httpx} & Solicitud del token al backend (\gls{tofu} y \gls{jwt}) \\
    \midrule
    \multicolumn{2}{@{}l}{\textit{Aplicación de usuario (Flutter, Dart)}} \\
    Dio & Cliente HTTP para consumir la \gls{api} \gls{rest} \\
    Riverpod & Gestión del estado de la aplicación \\
    \texttt{go\_router} & Navegación entre pantallas \\
    \texttt{shared\_preferences} & Persistencia local de la sesión \\
    \texttt{shadcn\_ui} & Componentes de la interfaz de usuario \\
    \texttt{url\_launcher} & Apertura del panel de Grafana en el navegador \\
    \texttt{intl} & Formato de fechas y localización \\
    \bottomrule
  \end{tabular}
\end{table}

\FloatBarrier

%% =============================================================================
\section{Criterios y métricas de evaluación}
\label{sec:metricas}

Dado que el sistema integra subsistemas de naturaleza distinta, su evaluación
contempla varias dimensiones complementarias:

\begin{itemize}
  \item \textbf{Navegación y patrulla.} Capacidad del robot para construir el
    mapa, generar una ruta de cobertura válida y recorrerla de forma cíclica sin
    colisiones, así como la tasa de éxito en alcanzar las poses objetivo.

  \item \textbf{Detección de incidencias.} Acierto del detector en la
    identificación de personas, considerando la confianza de las detecciones y
    la presencia de falsos positivos y falsos negativos.

  \item \textbf{Rendimiento del sistema de extremo a extremo.} Latencia desde
    que se produce un evento en el robot hasta que la incidencia queda
    registrada y es visible para el usuario, así como la regularidad del envío de
    telemetría.

  \item \textbf{Validación funcional de la integración.} Comprobación de que los
    flujos completos (arranque, mapeado, patrulla, detección, notificación y
    control desde la aplicación) se ejecutan de principio a fin.
\end{itemize}

Conviene señalar que, en el estado actual del trabajo, la evaluación realizada es
fundamentalmente de carácter \textbf{funcional}. La obtención de métricas
cuantitativas sistemáticas, especialmente sobre hardware real, queda pendiente y
se discute en el Capítulo~\ref{cap:resultados} y entre las líneas de trabajo
futuro del Capítulo~\ref{cap:conclusiones}.
